% Section: Budget Pacing Under Cost Drift
% Include from paper/sections/results.tex

\subsection{Budget Pacing Under Cost Drift (Figure~\ref{fig:exp03_adaptation})}
\label{sec:budget_drift}

\paragraph{Motivation.}
Experiments 01 and 02 studied budget pacing and non-stationarity in
isolation.  In production, both occur simultaneously: a router
operates under a cost budget when a model provider drops its pricing.
We ask: \emph{does the BudgetPacer automatically exploit a mid-stream
price drop to improve quality while maintaining budget compliance,
and how does this compare with unconstrained BanditGPT?}

\paragraph{Setup.}
The $K{=}3$ portfolio is deployed on a two-phase cost-drift stream:
Phase~1 ($n{=}893$) uses normal pricing (Gemini normalized cost
$\tilde{c}{=}0.67$); Phase~2 ($n{=}892$) applies a Gemini price
drop to \$0.10/\$0.10 per million tokens ($\tilde{c} \approx 0.0$).
Online learning uses the validation split ($n{=}1{,}785$) and
priors are trained on the disjoint training split
($n_{\mathrm{eff}}{=}10$, $\alpha{=}0.1$;
Section~\ref{sec:hparam_sweep}).

Three budget targets span the constraint regime:
\begin{itemize}[nosep,leftmargin=*]
  \item \textbf{Tight} ($B{=}\$2.3{\times}10^{-4}$/req):
    Lambda high, Llama-dominant in Phase~1.
  \item \textbf{Moderate} ($B{=}\$6.6{\times}10^{-4}$/req):
    Mixed Llama/Mistral routing in Phase~1.
  \item \textbf{Loose} ($B{=}\$1.9{\times}10^{-3}$/req):
    Light constraint in Phase~1.
\end{itemize}

\noindent
Per target, three pacer conditions are compared across 20~seeds:
\begin{enumerate}[nosep,leftmargin=*]
  \item \textbf{Pacer + $\gamma{=}0.997$}: Production config
    (BudgetPacer + jointly-tuned forgetting,
    effective memory ${\sim}333$ steps).
  \item \textbf{Pacer + $\gamma{=}1.0$}: Pacer-only (no forgetting).
  \item \textbf{Pacer + $\gamma{=}0.999$}: Pacer + mild forgetting
    (effective memory ${\sim}1{,}000$ steps).
\end{enumerate}

\noindent
An \textbf{unconstrained BanditGPT} baseline (same algorithm, same
priors, no budget pacer, $\lambda_s{=}0$) serves as quality ceiling.
Figure~\ref{fig:exp03_adaptation} uses the jointly-tuned config
(Pacer + $\gamma{=}0.997$) as the representative pacer condition;
variant ablation results are in Table~\ref{tab:exp03_ablation}.

\paragraph{Lambda adaptation (Figure~\ref{fig:exp03_adaptation}a).}
The dual variable $\lambda_t$ responds directly to the cost EMA.
At \textbf{moderate} budget, $\lambda_t$ oscillates around
$0.29$ during Phase~1 then drops sharply
to ${\sim}0$ within ${\sim}100$ steps of the price drop as
$\bar{c}_t$ falls below target.
At \textbf{loose} budget, $\lambda_t$ starts at $0.10$
and rapidly reaches zero after the price drop.
At \textbf{tight} budget, $\lambda_t$ decreases from $0.64$
to $0.39$ but stays elevated---the budget is so stringent that
even the EMA takes many observations to ratchet down from the
near-ceiling values accumulated during Phase~1.
This EMA-mediated response requires no external change-detection
signal: the dual update (Eq.~\ref{eq:dual_update}) naturally tracks
the cost regime.

\paragraph{Gemini adoption (Figure~\ref{fig:exp03_adaptation}b).}
At \textbf{moderate} budget, the pacer increases Gemini selection
from ${\sim}2\%$ to ${\sim}47\%$ within Phase~2, approaching the
unconstrained baseline (${\sim}68\%$ Gemini in Phase~2).
At \textbf{loose} budget, Gemini selection rises from ${\sim}10\%$
to ${\sim}46\%$, also converging toward the unconstrained level.
At \textbf{tight} budget, however, Gemini selection remains below
$1\%$ even after the price drop---an important limitation discussed
below.

The unconstrained BanditGPT routes $57\%$ to Gemini in Phase~1
(quality-optimal) and increases to $68\%$ after the price drop
removes the cost barrier entirely.

\paragraph{Cross-seed variance in routing paths.}
The $\pm$1\,SD bands in Figure~\ref{fig:exp03_adaptation}b are
wider for pacer conditions than for the unconstrained baseline,
particularly during the Phase~2 transition.
This is inherent to stochastic dual-variable optimization: different
prompt orderings produce different cost sequences, which drive
different $\lambda_t$ trajectories, which in turn produce different
arm-selection paths.  Three observations make this acceptable:
\begin{enumerate}[nosep,leftmargin=*]
  \item \textbf{Budget compliance is tight.}
    Despite the routing path variance, Phase~1 costs are within
    $1\%$ of target across all 20~seeds
    (Figure~\ref{fig:exp03_adaptation}c).  The constraint is
    satisfied even when the specific routing path varies.
  \item \textbf{Reward outcomes are stable.}
    The standard error on mean reward is ${\sim}0.001$--$0.002$
    across seeds (Table above), meaning quality outcomes are
    consistent even though routing paths differ.
  \item \textbf{Production sees one trajectory.}
    The bands show the distribution of possible paths across random
    prompt orderings, not noise within a single deployment.  Any
    individual deployment follows one smooth trajectory; the
    variance reflects sensitivity to input ordering, analogous to
    how stochastic gradient descent converges to similar losses from
    different initializations.
\end{enumerate}

\paragraph{Budget compliance (Figure~\ref{fig:exp03_adaptation}c).}

\begin{center}
\small
\begin{tabular}{@{}lrrrrr@{}}
\toprule
\textbf{Budget} & \textbf{Target} & \textbf{Phase~1} & \textbf{Phase~2} & \textbf{Overall} & \textbf{Util.} \\
\midrule
Tight    & \$2.34e-4 & \$2.37e-4 & \$2.33e-4 & \$2.35e-4 & $1.01\times$ \\
Moderate & \$6.62e-4 & \$6.64e-4 & \$3.49e-4 & \$5.07e-4 & $0.77\times$ \\
Loose    & \$1.87e-3 & \$1.83e-3 & \$3.52e-4 & \$1.09e-3 & $0.58\times$ \\
\midrule
Unconstr.\ & ---     & \$9.06e-3 & \$3.20e-4 & \$4.69e-3 & --- \\
\bottomrule
\end{tabular}
\end{center}

\noindent
All pacer conditions achieve Phase~1 costs within $1\%$ of the target.
After the price drop, costs decline below target (the pacer constrains
from above, never inflates); the overall utilization ratio therefore
falls below $1.0\times$.
The unconstrained baseline spends ${\sim}5$--$39\times$ more than pacer
targets in Phase~1 (\$9.1e-3/req), confirming that the pacer
provides genuine cost control.

\paragraph{Reward lift after price drop.}
We test whether the pacer exploits the price drop by measuring
$\Delta_\mathrm{reward} = \bar{r}_\mathrm{Phase\,2} -
\bar{r}_\mathrm{Phase\,1}$ within each condition (a paired
comparison; each condition serves as its own control).

\begin{center}
\small
\begin{tabular}{@{}llcccc@{}}
\toprule
\textbf{Budget} & \textbf{Condition} & \textbf{P1 Rwd} & \textbf{P2 Rwd} & $\boldsymbol{\Delta}$ & \textbf{95\% CI} \\
\midrule
Tight    & Pacer + $\gamma{=}0.997$ & 0.843 & 0.873 & $+0.030$ & $[0.025,\; 0.034]$ \\
Moderate & Pacer + $\gamma{=}0.997$ & 0.879 & 0.925 & $+0.046$ & $[0.042,\; 0.050]$ \\
Loose    & Pacer + $\gamma{=}0.997$ & 0.901 & 0.924 & $+0.023$ & $[0.019,\; 0.026]$ \\
\midrule
---      & Unconstrained            & 0.921 & 0.931 & $+0.010$ & $[0.009,\; 0.012]$ \\
\bottomrule
\end{tabular}
\end{center}

\noindent
All pacer conditions show highly significant reward improvement
($p < 10^{-11}$ for all; one-sample $t$-test on per-seed deltas,
19~d.f.).  The moderate target exhibits the largest lift
($\Delta{=}+0.046$), reflecting the biggest shift from
Llama-dominated to Gemini-heavy routing.
The unconstrained baseline shows a small but significant
\emph{increase} ($\Delta{=}+0.010$, $p < 10^{-10}$), reflecting the
fact that the price drop removes the cost barrier for Gemini, allowing
the router to exploit a quality-optimal model that was previously
penalized.

\paragraph{Pacer variant ablation.}

\begin{center}
\small
\label{tab:exp03_ablation}
\begin{tabular}{@{}llrrl@{}}
\toprule
\textbf{Budget} & \textbf{Condition} & \textbf{Regret} & $\boldsymbol{\Delta}$\textbf{\ vs.\ $\gamma$=0.999} & $\boldsymbol{p}$ \\
\midrule
\multirow{3}{*}{Tight}
  & $\gamma{=}0.997$ & $188.4 \pm 1.9$ & $-1.4$ & $0.536$ \\
  & $\gamma{=}1.0$   & $190.2 \pm 2.7$ & $+0.4$ & $0.848$ \\
  & $\gamma{=}0.999$ & $189.8 \pm 2.3$ & --- & --- \\[2pt]
\multirow{3}{*}{Moderate}
  & $\gamma{=}0.997$ & $109.8 \pm 1.8$ & $+6.5$ & $0.017^\dagger$ \\
  & $\gamma{=}1.0$   & $107.2 \pm 2.4$ & $+3.8$ & $0.209$ \\
  & $\gamma{=}0.999$ & $103.3 \pm 2.6$ & --- & --- \\[2pt]
\multirow{3}{*}{Loose}
  & $\gamma{=}0.997$ & $91.1 \pm 1.4$ & $\mathbf{+8.1}$ & $< 10^{-3}$ \\
  & $\gamma{=}1.0$   & $86.4 \pm 1.9$ & $+3.4$ & $0.048^\dagger$ \\
  & $\gamma{=}0.999$ & $83.0 \pm 1.0$ & --- & --- \\
\bottomrule
\end{tabular}
\end{center}

\noindent
Regret $\pm$1\,SE; $\Delta$ and $p$-values from two-sided paired
$t$-tests (19~d.f.).  Bold entries are significant after
Holm--Bonferroni correction at $\alpha{=}0.05$ across all six
comparisons; $\dagger$ denotes nominally significant ($p < 0.05$)
but not significant after correction.
Under cost shift, the ranking reverses compared to reward shift
(Experiment~02): $\gamma{=}0.999$ is consistently the \emph{best}
pacer condition, because reward estimates remain valid and any
forgetting needlessly decays them.
The jointly-tuned $\gamma{=}0.997$ is significantly worse at loose
budget ($+8.1$ regret, $p < 10^{-3}$), confirming that the
${\sim}333$-step effective memory optimized for reward shifts
forgets too aggressively when only costs change.
All conditions are indistinguishable at tight budget ($p \geq 0.5$),
where the dominant cost constraint overwhelms the forgetting effect.

\paragraph{Exploration starvation at tight budgets.}
At the tightest target, Gemini selection remains below 1\% even
after the price drop (Figure~\ref{fig:exp03_adaptation}b, orange).
Two factors combine:
\begin{enumerate}[nosep,leftmargin=*]
  \item \textbf{EMA momentum.}
    $\lambda_t$ was saturated near $0.64$ for much of
    Phase~1.  The smoothing constant $\alpha_\mathrm{ema}{=}0.05$
    creates ${\sim}14$-step inertia; many low-cost observations are
    needed to ratchet $\lambda_t$ down from its Phase~1 plateau.
  \item \textbf{Information deficit.}
    With $<1\%$ Gemini selection during Phase~1, the bandit's reward
    model for Gemini is almost entirely the warmup prior---it has not
    learned enough online to build a confident, calibrated estimate.
    Even when the cost penalty vanishes in Phase~2, the LinUCB
    exploration bonus is too small (with $n_\mathrm{eff}{=}10$
    prior) to pull Gemini above the well-explored Llama and Mistral
    arms.
\end{enumerate}

\noindent
This is a known property of budget-constrained bandits: very tight
budgets suppress exploration of expensive arms, and the resulting
information deficit persists even after the cost barrier is removed.
Practitioners should be aware that extremely tight budgets can create
persistent blind spots for arms that were historically cost-prohibitive.

\paragraph{Synthesis.}
This experiment demonstrates that the BudgetPacer automatically
exploits a mid-stream price drop to improve routing quality while
maintaining budget compliance---without any external intervention.
\begin{itemize}[nosep,leftmargin=*]
  \item The \textbf{pacer's dual-variable mechanism} tracks the
    new cost regime via the EMA, relaxing $\lambda_t$ and enabling
    Gemini routing at moderate and loose budgets
    (Figure~\ref{fig:exp03_adaptation}a--b).
  \item \textbf{Budget compliance is maintained} across both phases:
    Phase~1 costs are within 1\% of target, and the pacer naturally
    constrains from above (Figure~\ref{fig:exp03_adaptation}c).
  \item The \textbf{unconstrained baseline} already operates near the
    quality ceiling, spending ${\sim}5$--$39\times$ more.
    The pacer recovers most of the quality gap post-drop while
    maintaining strict cost control.
  \item Under pure cost shift, \textbf{less forgetting is better}:
    $\gamma{=}0.999$ outperforms $\gamma{=}0.997$ at loose budget
    ($p < 10^{-3}$) because reward estimates remain valid and
    need not be discounted.
    This contrasts with the reward shift in Experiment~02, where
    $\gamma{=}0.997$ is optimal---a tradeoff inherent to
    fixed-schedule forgetting.
  \item Very tight budgets create an \textbf{exploration starvation}
    effect: the bandit cannot learn about expensive arms it was
    prevented from selecting, and this information deficit persists
    after the cost barrier is removed.
\end{itemize}
